Para que la minería de datos se ejecute de forma estructurada, la industria utiliza metodologías estándar. La más famosa y utilizada a nivel mundial es \textbf{CRISP-DM} (\textit{Cross-Industry Standard Process for Data Mining}). Este marco de trabajo organiza todo proyecto en 6 fases cíclicas:

\begin{enumerate}
    \item \textbf{Comprensión del Negocio (Business Understanding):} Antes de procesar datos, se debe entender el problema a resolver. Se definen los objetivos comerciales y los criterios de éxito del proyecto.
    
    \item \textbf{Comprensión de los Datos (Data Understanding):} Se recolectan los datos iniciales y se realiza una exploración preliminar. Se identifican problemas de calidad (datos faltantes, errores) y se descubren pistas iniciales.
    
    \item \textbf{Preparación de los Datos (Data Preparation):} Suele ser la fase que más tiempo consume (70\% - 80\% del proyecto). Involucra limpiar bases de datos, transformar variables, unir tablas y formatear la información para los algoritmos.
    
    \item \textbf{Modelado (Modeling):} Se seleccionan y aplican diversas técnicas de Machine Learning (como las vistas anteriormente). Se calibran los parámetros (hiperparámetros) para encontrar el modelo con mejor rendimiento.
    
    \item \textbf{Evaluación (Evaluation):} Se analiza el modelo rigurosamente. No solo se verifica su precisión matemática, sino que se valida si realmente cumple con el objetivo de negocio planteado en la primera fase.
    
    \item \textbf{Despliegue (Deployment):} El modelo se lleva al entorno real. Puede implementarse como un reporte técnico, un panel interactivo (\textit{dashboard}) o integrarse como un sistema automatizado.
\end{enumerate}